\section{Data}

\subsection{Interaction data}

The measured data consist of positive observations only. They come from a system-wide lipidomics
screen of lipid transfer proteins (LTPs) \cite{acg}. Each protein was assayed in two screens
against a set of lipid species, and the species recovered from a protein were recorded as bound.
The screen reports 634 such (protein, lipid) combinations. Each carries the screen or screens the
species was recovered in, and 50 were recovered in both.

The panel comprises 35 proteins in nine families and 283 lipid species in 34 head-group classes
(Table~\ref{tab:families}).

\begin{table}[htbp]
\centering
\begin{tabular}{lrrp{0.52\linewidth}}
\hline
family & prot. & pos. & positives per protein \\
\hline
CRAL-TRIO & 9 & 165 & RLBP1 62, TTPAL 62, SEC14L6 14, TTPA 10, SEC14L2 9, SEC14L4 3, ATCAY 2, SEC14L5 2, BNIPL 1 \\
START & 3 & 150 & STARD10 65, STARD2 64, STARD11 21 \\
lipocalin & 10 & 87 & LCN1 27, FABP1 24, FABP5 10, PMP2 9, FABP7 8, LCN15 5, CRABP2 1, RBP1 1, RBP4 1, RBP5 1 \\
GLTP & 2 & 62 & GLTP 33, GLTPD1 29 \\
IP\_trans & 3 & 58 & PITPNB 26, PITPNA 21, PITPNC1 11 \\
LBP\_BPI\_CETP & 2 & 53 & BPI 29, BPIFB2 24 \\
scp2 & 3 & 43 & SCP2D1 28, SCP2 10, HSDL2 5 \\
ML & 1 & 10 & GM2A 10 \\
OSBP & 2 & 6 & OSBPL5 3, OSBPL9 3 \\
\hline
\end{tabular}
\caption{The panel: proteins per family, and positive combinations per family and per protein.}
\label{tab:families}
\end{table}

The negative class is built as the complement. The panel defines $35 \times 283 = 9905$
combinations, and the 9271 that were not reported bound receive the label zero. The training
table is therefore the complete grid, one row per combination, at a positive rate of 6.4\%. Its
zeros are unlabelled combinations, not measured non-binding, and an unknown fraction of them are
species that bind but escaped detection. That rate is a property of the completion rule as much as
of the biology.

\subsection{Protein structures}

Each protein was assigned one structure. Twenty-four are experimental structures from the PDB: 22
by X-ray diffraction, one by solution NMR and one by joint X-ray and neutron diffraction. The
remaining 11 are AlphaFold DB models \cite{alphafold}, used where no suitable experimental
structure was available. Candidates were selected on Uniprot family and domain annotation, largely
from PROSITE, and structures carrying mutations were excluded.

Deposited files were normalised before annotation. Selenomethionine is stored as a heteroatom
record, so the default protein filter removed it and silently shortened two chains. Rewriting it
as methionine beforehand gives GM2A and PITPNA 162 and 269 residues, in agreement with their
sequences. Residues unresolved in the crystal are absent from the coordinate file and are left
absent, without placeholders, so a contact may form between residues distant in sequence; CERT has
two such gaps, at 493--495 and 535--540. No residue is discarded on predicted confidence and no
domain is cropped. The chain enters the representation as deposited.

\subsection{Pocket annotation}

Cavities were annotated with the pocket detection of Voronota \cite{voronota}, which operates on
the Voronoi tessellation of the atomic balls. The same tessellation yields the residue contacts
and the burial values used below. Cavity membership, neighbourhood and burial therefore derive
from one geometric model and refer to the same residues. LTP cavities are large and hydrophobic,
and the annotation has to separate the interior cavity from surface grooves, so detection
parameters were tuned per protein and every annotated structure was inspected visually.
Alternative cavity detection algorithms were not benchmarked. The choice rests on the consistency
of the representation and on the visual check, not on a measured advantage over them.

A residue is assigned to the cavity when any of its side-chain atoms is annotated. Backbone atoms
and the beta carbon are excluded, so membership reflects the side chains lining the cavity. Burial
is kept as a continuous per-residue quantity rather than a second binary mask. Its defined values
span 0.005 to 0.989 across the 35 structures. The software emits the value 2 where burial is
undefined; the 716 residues of 7804 (9.2\%) that carry it are passed on unchanged.

\subsection{Structure graphs}

Each structure is a graph over its residues. Two residues are joined when their Voronoi cells
share a face, so adjacency follows contact rather than a distance threshold or a fixed number of
nearest neighbours. Nodes carry residue type, solvent-accessible surface area and residue volume,
together with burial, cavity membership and the language-model embedding of the next subsection.
Edges carry the distance between the residues, the area of their contact and the solvent-facing
part of that contact.

The same construction yields an atom-level graph. The residue-level one is used for two reasons.
With 35 structures the atom-level description is too fine to fit without memorisation. And the
residue level coincides with the resolution of the protein language model, which allows the
geometric and the evolutionary description to be concatenated node by node.

\subsection{Protein embeddings}

Geometric and chemical descriptors alone proved insufficient: no configuration trained on them
exceeded the trivial predictor. The graph is therefore conditioned on a protein language model.
Sequences are embedded with ESM3 \cite{esm3} and the residue embeddings are attached to the
corresponding nodes. An embedding of $N$ residues has $N+2$ rows, whose terminal tokens are
removed, so that nodes, embedding rows and annotated residues agree in number and order. The
structural normalisation above exists to protect this correspondence, which a residue missing from
the coordinate file but present in the sequence breaks silently. And 1536 embedding dimensions
beside three structural attributes would dominate the node representation, so the embedding is
projected linearly to a low dimension before concatenation.

These embeddings are computed from sequence alone, structural information reaching the model
through the graph.

What these embeddings distinguish was measured directly. Averaged over residues, the 35 proteins
are almost one point: over the 595 pairs the median cosine between two ESM3 vectors is 0.974, with
a floor of 0.905. Their binding profiles are the opposite. The median cosine between two profiles
is 0.000, and 59\% of pairs share no class at all. The embedding therefore separates proteins by
family membership, not by what they bind: a protein's nearest neighbour in ESM3 space belongs to
its own family 94\% of the time, against 16\% for a random neighbour.

Table~\ref{tab:proteinrepr} asks whether any representation beats ignoring the protein. None does.
ESM-IF1 \cite{esmif}, an inverse-folding encoder reading backbone coordinates, spreads the cloud
out and still loses, so the tight cloud is a symptom rather than the cause. ProteinMPNN
\cite{proteinmpnn} packs the proteins tighter than ESM3 and lands between the two. ESM3 itself
beats the reference on average, but on four of the seven families its neighbours are worse than
ignoring the protein. The protein side is limited by 35 objects, not by the representation chosen
for them.

Hand-built cavity descriptors were measured under an easier regime, with a single protein held out
rather than a whole family, so their relatives stay in training. There ESM3 reconstructs the
profile at 0.331 against 0.249 for ignoring the protein, thirteen cavity descriptors reach 0.274
and the four shape numbers among them -- lining volume over open surface, cavity extent,
elongation and flatness -- 0.228, correlating with the profile at $+0.096$ and $+0.110$. Appending
the thirteen to five principal components of ESM3 gives 0.294 against 0.279 for those components
alone, so they add a little there and still fall short of the whole embedding. Their
nearest-neighbour family rate is 0.51 and 0.63
(\texttt{preprocessing/pocket\_descriptor\_identity\_check.py}).

\begin{table}[htbp]
\centering
\begin{tabular}{lrrrrrr}
\hline
representation & dim & NN same family & median cos & min cos & corr.\ profile & profile, family out \\
\hline
ignore the protein & -- & -- & -- & -- & -- & 0.151 \\
ESM3 & 1536 & 0.943 & 0.974 & 0.905 & $+0.235$ & \textbf{0.175} \\
ESM-IF1 & 512 & 0.800 & 0.762 & 0.256 & $+0.134$ & 0.137 \\
ProteinMPNN & 128 & 0.857 & 0.994 & 0.958 & $+0.107$ & 0.124 \\
\hline
\end{tabular}
\caption{What each protein representation is worth, on the 35 mean-pooled vectors the model would
receive. \emph{NN same family}: how often a protein's nearest other protein in that representation
shares its family, against 0.156 for a random neighbour. \emph{median cos} and \emph{min cos}: how
far apart the representation puts two proteins at all. \emph{corr.\ profile}: correlation between
representation similarity and binding-profile similarity over the 595 pairs. \emph{profile, family
out}: cosine between a held-out protein's true binding profile and the one reconstructed from its
three nearest training proteins, averaged over the seven families. Computed by
\texttt{preprocessing/protein\_representation\_identity\_check.py}.}
\label{tab:proteinrepr}
\end{table}

\subsection{Lipid structures}

The screen identifies a species by head group, total acyl carbon count and total number of double
bonds, and structures consistent with these identities were retrieved from LIPID MAPS
\cite{lipidmaps} in SMILES form.

Mass spectrometry constrains the head group and the total mass. It frequently does not constrain
the distribution of carbons between the chains, the position of the unsaturations, or which chain
occupies the \emph{sn}-1 and which the \emph{sn}-2 position. This ambiguity is represented
explicitly: a species is stored as a set of candidate structures, with duplicates merged after
canonicalisation. Of the 9905 rows, 6999 (70.7\%) carry more than one candidate. Phosphatidyl\-
ethanolamine (34:1), for example, is represented by four 16:0/18:1 structures, differing in the
position of the double bond and in the order of the chains on the glycerol backbone.

A row is encoded by one of its candidates, drawn afresh each time the row is presented during
training. Over an epoch a species is therefore shown as its whole candidate set, not as one
arbitrary member of it. Validation and test read the first-listed candidate instead and keep it
fixed across epochs, so the metric that selects the checkpoint and stops training early moves for
the model rather than for which structure a row happened to show it.

The retrieved strings carry the chirality of the glycerol carbon and the geometry of the double
bonds. The encoding table used by default, however, is built after canonicalisation without
stereochemical descriptors, so none of the 1226 distinct structures reaches the chemical language
model with its stereochemistry intact. An isomeric table is available, as is a molecular-graph
representation carrying chirality and bond geometry as explicit atom and bond attributes. Whether
the chemical language model separates stereoisomers when they are supplied has not been tested
here.

\subsection{Lipid embeddings}

Lipids are encoded with MolFormer \cite{molformer}, a transformer pretrained at scale on SMILES
strings. It is used frozen: the encodings of all distinct candidate structures are computed once
and no gradient reaches the encoder. A chemical language model matches the level of description
available for the lipid, a formula and a connectivity rather than a conformation. Its token-level
output also lets the protein side attend over the structure instead of receiving one pooled
vector.

Canonical strings are tokenised with the model's own tokeniser and passed through the pretrained
encoder, and the token-level output is retained. A structure of $T$ tokens is stored as $T+2$
vectors of dimension 768. The two additional positions are the terminal tokens introduced by the
tokeniser, the leading one being the sequence-level token the model was trained to summarise a
structure with. Token counts range from 22 to 130 over the 1226 structures, 24 to 132 rows stored.

One property of the architecture affects reproducibility. The linear attention of MolFormer draws
a random feature matrix at each forward pass. Our first table was computed without fixing that
draw, and two encodings of the same structure differed there by as much as 8.8 on a scale of $\pm
10$. The table used here was computed with the matrix stored in the checkpoint, so figures
obtained with the earlier table are not comparable with those reported below.

The alternative lipid representation is an explicit atom-bond graph with eleven atom and six bond
attributes. Molecular fingerprints enter only as a similarity measure between lipids, in the
sampling weights and in the chemistry baselines, never as model input.

\subsection{Data splits}
\label{sec:splits}

Because the matrix is filled completely, holding out a set of rows on one axis leaves the other
axis intact. A species from a held-out protein family still appears in 34 further rows, so its
frequency of binding across the panel is known. That frequency predicts the held-out labels
without any information about the held-out protein. Every split is therefore reported against a
reference predictor that ignores the protein and answers with the training frequency of the
species, or of its head-group class. Computed on the rows the model receives, it measures how much
of the metric a lookup can produce. All splits are group-disjoint; no rows are split at random.

The model does not see the whole grid. Every positive is kept, and negatives are drawn to match,
two per positive within each protein (\texttt{balanced\_proteins}); these rows are the pool. That
leaves about one third of it positive and equalises the positive rate across proteins, so no
protein can be answered by its own base rate. The rows a split holds out of the pool are its
block, halved into validation and test. A seed changes two things and nothing else: which zeros
are drawn, and how the block is halved. Which proteins and which lipid classes leave training is
fixed. Every figure below is averaged over five seeds.

The same ratio holds one level down, inside a batch: \texttt{balanced\_batches} draws positives
and unlabelled rows two to one rather than splitting a batch evenly between them. Each class is
chunked into as many pieces as an epoch needs batches, sizes differing by at most one row, so
every row of the pool is seen exactly once per epoch and no batch is left carrying only one
class. Matching negatives to positives once across the whole table instead of per protein, and
splitting a batch evenly between classes instead of by the pool's own ratio, are both available
and were not used; \texttt{balanced\_proteins} and \texttt{balanced\_batches} are on together in
every configuration reported below.

\paragraph{Holding out proteins.} One or more families leave training and their rows are divided
equally between validation and test. Alternatively, validation and test are two different
families. In the pool, 99.9\% of test species and all head-group classes have been seen in
training. The reference predictor accordingly reaches 0.582 balanced accuracy on test averaged
over the nine families, and up to 0.68 on individual ones. No training row repeats as a validation
row, so the pairs are new: what carries over is the frequency of the species, not the pair.
Results obtained under this split are to be read against 0.58, not 0.50.

\paragraph{Holding out lipid chemistry.} The complementary split keeps every protein in training
and removes a chemical family of lipids. Four sets were defined by chemistry: a set is only unseen
if its close relatives leave with it. Similarity to the training chemistry measures how unseen a
set is: the mean, over the structures of the set, of the highest fingerprint similarity to any
structure remaining in training (Table~\ref{tab:lipidsets}). Lower is more isolated. The first
three sets are well separated. The anionic set is not, and cannot be: the anionic
glycerophospholipids differ from the phosphatidylcholines that remain only in the head group,
while a fingerprint is dominated by the two acyl chains. It is retained as the most demanding of
them because it asks whether the head group is read at all.

\begin{table}[htbp]
\centering
\begin{tabular}{llrl}
\hline
set & similarity & positives held out & content \\
\hline
sphingolipids & 0.458 & 66 (10.4\%) & sphingoid backbone, all of it \\
phosphorus-free & 0.553 & 61 (9.6\%) & neutral glycerolipids, free fatty acids, retinol \\
choline & 0.653 & 223 (35.2\%) & phosphocholine head, di- and lyso- \\
anionic & 0.766 & 186 (29.3\%) & anionic glycerophospholipid heads \\
\hline
\end{tabular}
\caption{Lipid sets held out of training, ordered by their similarity to the chemistry that stays
in training. Lower is more isolated.}
\label{tab:lipidsets}
\end{table}

\paragraph{Holding out both axes.} The split used below cuts both axes at once. A protein family
leaves training, and a set of head-group classes leaves training for every protein, not for that
family alone. A row of the family in a held-out class then has neither its protein nor its lipid
chemistry in training, and no frequency can be estimated for it.

Which classes leave is decided per family. A family's positives sit in the classes it transports,
so a fixed partition would be arbitrary for whichever family is held out. Each class $c$ is scored
by
\[
s(c) = \frac{\text{positives of the held-out family in } c}{\text{positives of all other proteins in } c + 1},
\]
high for a class the family owns, low for one the rest of the panel also uses. Classes are removed
in decreasing order of $s$ until they cover a set share of the family's positives. Two kinds of
row are then discarded rather than evaluated. The family's rows in the classes that stayed go
first: their protein is absent, so they cannot train, and they are the only rows a frequency
lookup could still score. The other proteins' rows in the removed classes go with them. Their
chemistry is unseen but their own family sits in training, so scoring them answers the one-axis
question rather than this one, and there are enough of them to decide the answer: left in, a
held-out family held between 21 and 44 per cent of its own block for six of the seven families,
GLTP the exception at 97 per cent, since almost nothing else in the panel binds a sphingolipid.
Together the two discards cost between 2 and 35 per cent of the working set, and nothing from
training, which never held either.

That discarding fixes the reference predictor at 0.500 for every family and every threshold. The
threshold lowers that similarity instead. The classes the selection stops short of are the nearest
relatives of the family's own lipids -- for the phosphatidylcholines of START, the
lysophosphatidylcholines -- and while they stay in training every held-out structure keeps a close
neighbour to extrapolate from. Table~\ref{tab:coldsplitshare} gives that similarity per family and
per threshold, measured as above (\texttt{analysis/coldsplit\_geometry.py}).

\begin{table}[htbp]
\centering
\begin{tabular}{lrrrrrrr}
\hline
held-out family & 0.30 & 0.50 & 0.70 & 0.75 & \textbf{0.80} & 0.85 & 0.90 \\
\hline
CRAL-TRIO & 0.838 & 0.796 & 0.796 & 0.796 & \textbf{0.783} & 0.765 & 0.779 \\
GLTP & 0.728 & 0.697 & 0.697 & 0.697 & \textbf{0.567} & 0.571 & 0.599 \\
IP\_trans & 0.793 & 0.793 & 0.793 & 0.793 & \textbf{0.793} & 0.803 & 0.803 \\
LBP\_BPI\_CETP & 0.793 & 0.793 & 0.793 & 0.793 & \textbf{0.793} & 0.793 & 0.795 \\
lipocalin & 0.890 & 0.775 & 0.795 & 0.795 & \textbf{0.795} & 0.805 & 0.805 \\
scp2 & 0.797 & 0.854 & 0.854 & 0.854 & \textbf{0.854} & 0.831 & 0.831 \\
START & 0.792 & 0.792 & 0.792 & 0.794 & \textbf{0.794} & 0.794 & 0.804 \\
\hline
mean & 0.804 & 0.786 & 0.789 & 0.789 & \textbf{0.768} & 0.766 & 0.774 \\
positives left in training, lowest & 353 & 353 & 327 & 293 & \textbf{293} & 225 & 116 \\
\hline
\end{tabular}
\caption{Similarity of the held-out block to the training chemistry, per family and per threshold:
the mean over the block's structures of the highest Tanimoto similarity to a structure still in
training. Lower is more isolated. The column in use is 0.80.}
\label{tab:coldsplitshare}
\end{table}

The similarity does not fall monotonically with the threshold, since a newly removed class can
bring in structures whose own relatives stay behind. The value used is 0.80. It reaches a mean
similarity of 0.768 and still leaves every family at least 293 positives to train on; beyond it
the gain is small and the cost is not. One threshold serves all families even though their optima
differ. scp2 is the one to read with care: it stays near 0.85 wherever the threshold is put, so its
block is never a test of unseen chemistry.

Table~\ref{tab:doublesplit} reports the resulting blocks. ML and OSBP cannot be run: with ten and
six positives in the whole family, no block can be assembled at any threshold.

\begin{table}[htbp]
\centering
\begin{tabular}{lrrrrrr}
\hline
held-out family & classes removed & positives in block & train positives lost & train & valid & test \\
\hline
CRAL-TRIO & 7 & 134 & 119 & 1044 & 129 & 128 \\
GLTP & 5 & 51 & 1 & 1713 & 52 & 50 \\
IP\_trans & 3 & 47 & 213 & 1089 & 71 & 70 \\
LBP\_BPI\_CETP & 4 & 47 & 193 & 1164 & 71 & 70 \\
lipocalin & 14 & 72 & 220 & 981 & 108 & 108 \\
scp2 & 7 & 34 & 126 & 1395 & 51 & 51 \\
START & 4 & 129 & 191 & 879 & 153 & 154 \\
\hline
\end{tabular}
\caption{Blocks produced by the two-axis split, at threshold 0.80 and two negatives per positive.
Every count is fixed by the split rule rather than averaged: the seed changes which zeros are drawn
and how the block is halved, not how many rows fall where, so these figures are identical across the
five seeds.}
\label{tab:doublesplit}
\end{table}

Two properties of these blocks matter when per-family results are read. The five classes GLTP
loses -- sphingomyelin, hexosylceramide, dihexosylceramide, sulfohexosylceramide and ceramide
phosphate -- hold exactly one positive anywhere else in the panel, the sphingomyelin LCN1 takes.
Removing them costs training that single row, and the model therefore reaches the GLTP block
having seen no positive sphingolipid at all. That block is also the smallest, at 102 rows against
START's 307 and tied with scp2, so GLTP and scp2 are the two families whose numbers carry the
least weight. Second, the positive rate of a block is not the training rate. Negatives are matched
two per positive, but only as far as the removed classes hold any: a family is held out on the
classes it owns, and inside those classes most of its cells are positive. The blocks accordingly
run from 33 per cent positive, where the match is met in full, to 50 for GLTP and 52 for
CRAL-TRIO, where it is not. Balanced accuracy is indifferent to that; average precision compared
across families is not.

The two operations have to be ordered. Removing classes cuts across proteins, so a protein whose
positives lay in a removed class loses nearly all of its positives and almost none of its
negatives. Sampling negatives over the whole table and cutting afterwards left LCN1 in training
with one positive against 45, and put as many as eight proteins per held-out family into training
carrying a single label value, which no reweighting can repair. Negatives are therefore matched to
positives separately within each side of the coming split. Over the seven families and five seeds
no protein then reaches training one-sided, and the matching discards nothing of its own: the rows
that go are the two kinds named above.

\paragraph{What the splits guarantee.} A split is accepted only once the reference predictor
reaches chance with no species in common between training and test (Table~\ref{tab:baselines}).
What it still does not remove is extrapolation by chemical similarity. An unseen species can
resemble one that was seen, and resemblance to a training positive is again a prediction that
requires no protein. That residual is measured rather than eliminated, which is why the reference
for ranking quality is a chemistry-based null model rather than chance.

\begin{table}[htbp]
\centering
\begin{tabular}{lrr}
\hline
split & reference predictor (test) & test species seen in training \\
\hline
protein families held out & 0.582 & 99.9\% \\
two axes, familiar rows retained & 0.508 & 7\% \\
two axes & 0.500 & 0\% \\
lipid chemistry held out & 0.500 & 0\% \\
\hline
\end{tabular}
\caption{Balanced accuracy of a predictor that ignores the protein and answers with the training
frequency of the species, and the corresponding overlap of species between training and test, over
every family that can be held out: nine for the protein split, seven for the two-axis one.}
\label{tab:baselines}
\end{table}

With nine families of very unequal size no single split is representative of the panel, so each
family that can be held out is held out in turn and reported on its own.